% from ega1-10-fr.tex
\oldpage[I]{183}
(resp. $\mathfrak{j}_{\psi(x)}$)~; la remarque précédente montre donc que
${}^a\varphi=\psi$. Enfin, pour tout $g\in B$, le diagramme
\[
  \xymatrix{
    B=\Gamma(\mathfrak{Y},\mathscr{O}_{\mathfrak{Y}})
      \ar[r]^\varphi\ar[d] &
    \Gamma(\mathfrak{X},\mathscr{O}_{\mathfrak{X}})=A\ar[d]\\
    B_{\{g\}}=\Gamma(\mathfrak{D}(g),\mathscr{O}_{\mathfrak{Y}})
      \ar[r]_{\Gamma(\theta_{\mathfrak{D}(g)})} &
    \Gamma(\mathfrak{D}(\varphi(g)),\mathscr{O}_{\mathfrak{X}})
      =A_{\{\varphi(g)\}}
  }
\]
est commutatif~; d'après la propriété universelle des anneaux complets de
fractions (\hyperref[0.7.6.6-fr]{0, 7.6.6}),
$\theta_{\mathfrak{D}(g)}$ est égale à
$\widetilde{\varphi}_{\mathfrak{D}(g)}$ pour tout $g\in B$, donc
(\hyperref[0.3.2.5-fr]{0, 3.2.5}) on a
$\theta=\widetilde{\varphi}$.
\end{proof}

Nous dirons qu'un morphisme $(\psi,\theta)$ d'espaces topologiquement
annelés vérifiant la condition de (\hyperref[I.10.2.2-fr]{10.2.2}) est un
\emph{morphisme de schémas formels affines}. On peut dire que les foncteurs
$\operatorname{Spf}(A)$ en $A$ et
$\Gamma(\mathfrak{X},\mathscr{O}_{\mathfrak{X}})$ en $\mathfrak{X}$
définissent une \emph{équivalence} de la catégorie des anneaux admissibles
et de la catégorie duale de la catégorie des schémas formels affines
(T, I, 1.2).

\begin{env}[10.2.3]
\label{I.10.2.3-fr}
Comme cas particulier de (\hyperref[I.10.2.2-fr]{10.2.2}), notons que, pour
$f\in A$, l'injection canonique du schéma formel affine induit par
$\mathfrak{X}$ sur $\mathfrak{D}(f)$ correspond à l'homomorphisme continu
canonique $A\to A_{\{f\}}$. Sous les hypothèses de
(\hyperref[I.10.2.2-fr]{10.2.2}), soient $h$ un élément de $B$, $g$ un
élément de $A$, multiple de $\varphi(h)$~; on a alors
$\psi(\mathfrak{D}(g))\subset\mathfrak{D}(h)$~; la restriction de $u$ à
$\mathfrak{D}(g)$, considérée comme morphisme de $\mathfrak{D}(g)$ dans
$\mathfrak{D}(h)$, est l'unique morphisme $v$ rendant commutatif le
diagramme
\[
  \xymatrix{
    \mathfrak{D}(g)\ar[r]^v\ar[d] &
    \mathfrak{D}(h)\ar[d]\\
    \mathfrak{X}\ar[r]_u &
    \mathfrak{Y}
  }
\]
Ce morphisme correspond à l'unique homomorphisme continu
$\varphi':B_{\{h\}}\to A_{\{g\}}$
(\hyperref[0.7.6.7-fr]{0, 7.6.7}) rendant commutatif le diagramme
\[
  \xymatrix{
    A\ar[d] &
    B\ar[l]_\varphi\ar[d]\\
    A_{\{g\}} &
    B_{\{h\}}\ar[l]_{\varphi'}
  }
\]
\end{env}

\subsection{Idéaux de définition d'un schéma formel affine.}
\label{subsection:I.10.3-fr}

\begin{env}[10.3.1]
\label{I.10.3.1-fr}
Soient $A$ un anneau admissible, $\mathfrak{J}$ un idéal ouvert de $A$,
$\mathfrak{X}$ le schéma formel affine $\operatorname{Spf}(A)$. Soit
$(\mathfrak{J}_\lambda)$ l'ensemble des idéaux de définition de $A$
contenus dans $\mathfrak{J}$~; alors
$\widetilde{\mathfrak{J}}/\widetilde{\mathfrak{J}}_\lambda$ est un
faisceau d'idéaux de
$\widetilde{A}/\widetilde{\mathfrak{J}}_\lambda$. Désignons par
$\mathfrak{J}^\Delta$ la limite projective des faisceaux induits sur
$\mathfrak{X}$ par
$\widetilde{\mathfrak{J}}/\widetilde{\mathfrak{J}}_\lambda$, qui
s'identifie à un \emph{faisceau d'idéaux} de
$\mathscr{O}_{\mathfrak{X}}$ (\hyperref[0.3.2.6-fr]{0, 3.2.6}). Pour tout
$f\in A$, $\Gamma(\mathfrak{D}(f),\mathfrak{J}^\Delta)$ est la limite
projective de
$S_f^{-1}\mathfrak{J}/S_f^{-1}\mathfrak{J}_\lambda$, autrement dit
s'identifie à l'idéal ouvert $\mathfrak{J}_{\{f\}}$ de l'anneau
$A_{\{f\}}$ (\hyperref[0.7.6.9-fr]{0, 7.6.9}), et en particulier
$\Gamma(\mathfrak{X},\mathfrak{J}^\Delta)=\mathfrak{J}$~; on en conclut
(les $\mathfrak{D}(f)$ formant une base de la topologie de $\mathfrak{X}$)
que l'on a
\[
  \mathfrak{J}^\Delta|\mathfrak{D}(f)
  =(\mathfrak{J}_{\{f\}})^\Delta.
  \tag{10.3.1.1}\label{I.10.3.1.1-fr}
\]
\end{env}

